\documentclass[11pt]{article}

\usepackage{hw_style}

% Set due date for this homework (updated to 2025)
\setduedate{10}{29}

% Setup homework number and header
\hwnumber{5}

\begin{document}

% \begin{hwproblem}{1}
%   计算积分
%   $$
%   \int_0^{2 \pi} \cos ^{2 n} x d x \text {. }
%   $$
%   \end{hwproblem}

% \vspace{0.5cm}
  
% \begin{hwproblem}{2}
%   计算积分
%   $$
%   \int_{-\infty}^{\infty} \frac{e^{\ii m x}}{x-\ii \alpha} d x,
%   $$
%   和
%   $$
%   \int_{-\infty}^{\infty} \frac{e^{\ii m x}}{x+\ii \alpha} d x.
%   $$
%   其中$m>0, \Re \alpha > 0$.
% \end{hwproblem}

\vspace{0.5cm}

\begin{hwproblem}{}
  试推导多目标变量泛函$J\left(u_1, u_2, \cdots, u_n, u_1^{\prime}, u_2^{\prime}, \cdots, u_n^{\prime} \mid x\right)$
  取极值的欧拉方程是
  $$
  \frac{\partial L}{\partial u_i}-\frac{\partial}{\partial x}\left(\frac{\partial L}{\partial u_i^{\prime}}\right)=0 .
  $$
  \end{hwproblem}

\vspace{0.5cm}
  
  
\begin{hwproblem}{}
  过二已知点 $\left(x_1, y_1\right),\left(x_2, y_2\right)$ 作一曲线, 使此曲线绕 $x$ 轴旋转所得曲面面积最小, 求曲线满足的微分方程.
\end{hwproblem}

\vspace{0.5cm}
    

\begin{hwproblem}{}
求泛函
$$
J[y] = \int_0^1 \left[ \left(x \frac{dy}{dx} \right)^2 + 2 y^2 \right] dx
$$
在边界条件为$y(0) = 0, y(1) = 1$下的最小值.
\end{hwproblem}

\vspace{0.5cm}

\begin{hwproblem}{}
  请求解最速降线的解, 它对应的泛函表示为
  $$
  T[u] = \int_{x_0}^{x_1} \frac{\sqrt{1+u^{\prime 2}}}{\sqrt{2 g\left(y_0-u\right)}} d x .
  $$
\end{hwproblem}

\vspace{0.5cm}

% \begin{hwproblem}{7}
% 证明对于矩形波,
% $f(x)=\left\{ \begin{array}{ll}-1, & x<0 \\ +1, &  x>0\end{array}\right.$, 其傅里叶级数为
% $$
% f(x)=\frac{4}{\pi} \left[  \sin \pi x+\frac{1}{3} \sin 3 \pi+\frac{1}{5 } \sin 5 x+\cdots \right].
% $$
% \end{hwproblem}

% \vspace{0.5cm}

\hwrequirements

\end{document}